% Generated draft: requires mathematical and source review.
\documentclass[11pt]{article}
\usepackage[margin=25mm]{geometry}
\usepackage{fontspec}
\setmainfont{Latin Modern Roman}
\usepackage{amsmath,amssymb,mathtools,unicode-math}
\setmathfont{Latin Modern Math}
\usepackage{enumitem,needspace,xurl,hyperref}
\setlength{\parindent}{0pt}
\setlength{\parskip}{5pt}
\setlength{\emergencystretch}{3em}
\newcommand{\N}{\mathbb{N}}
\newcommand{\Z}{\mathbb{Z}}
\newcommand{\Q}{\mathbb{Q}}
\newcommand{\R}{\mathbb{R}}
\newcommand{\C}{\mathbb{C}}
\newcommand{\F}{\mathbb{F}}
\newcommand{\A}{\mathbb{A}}
\newcommand{\PP}{\mathbb{P}}
\newcommand{\E}{\mathbb{E}}
\newcommand{\eps}{\varepsilon}
\newcommand{\dd}{\,\mathrm{d}}
\begin{document}
\section*{UnsolvedMath ID: 30005694}
\textbf{Source problem code:} catalogue entry
\subsection{Definitions and mathematical statement}
A totally disconnected locally compact group is a locally compact Hausdorff topological group whose connected components are singletons. For a topological group \(G\), write \(\Q[G]\text{-}\mathrm{dis}\) for the category of rational vector spaces equipped with a left \(G\)-action such that every vector has an open stabilizer, with \(G\)-equivariant \(\Q\)-linear maps as morphisms. An object \(P\) is projective if, for every surjective morphism \(f\colon X\to Y\) and every morphism \(u\colon P\to Y\), there is a morphism \(v\colon P\to X\) satisfying \(f\circ v=u\).

\textbf{Tensor-projectivity question.} In the intended setting of the source, let \(P\) be any projective object of \(\Q[G]\text{-}\mathrm{dis}\), and let \(M\) be any object of that category. Equip \(P\otimes_{\Q}M\) with the diagonal action
\[
 g\cdot(p\otimes m)=(g\cdot p)\otimes(g\cdot m).
\]
Must \(P\otimes_{\Q}M\) be projective in \(\Q[G]\text{-}\mathrm{dis}\) for every such pair \((P,M)\)? The diagonal action is discrete because the stabilizer of a finite sum of pure tensors contains a finite intersection of open stabilizers. The workshop context indicates totally disconnected locally compact groups, but the supplied formulation does not explicitly specify the class of groups or any additional hypotheses on \(G\) for this first question; its group quantifier therefore remains unresolved.

\textbf{Stallings--Swan extension question.} For a totally disconnected locally compact group \(G\), let
\[
 \operatorname{cd}_{\Q}(G)
 =\operatorname{pd}_{\Q[G]\text{-}\mathrm{dis}}(\Q),
\]
where \(\Q\) carries the trivial action. A group is compactly generated if a compact subset generates it as a group, and it is unimodular if its modular function is identically \(1\). Retaining compact generation and the hypothesis \(\operatorname{cd}_{\Q}(G)\leq 1\), the source asks whether the conclusion of its \emph{Theorem 2}, described as a Stallings--Swan theorem, remains valid upon removing compact-open boundedness alone, unimodularity alone, or both hypotheses together. The supplied record does not reproduce the conclusion of Theorem 2 or define compact-open boundedness. Consequently, this second question specifies the requested changes of hypotheses but cannot be rendered as a complete mathematical implication from the supplied evidence.
\subsection{Short English statement}
The first question asks whether tensoring a projective rational discrete \(G\)-module with any rational discrete \(G\)-module preserves projectivity under the diagonal action. The second asks whether a cited Stallings--Swan theorem extends after dropping boundedness, unimodularity, or both, while retaining compact generation and rational cohomological dimension at most one. The first question's exact group hypotheses and the second question's theorem conclusion and boundedness definition require clarification.
\subsection{Sources}
\begin{itemize}
\item Local UnsolvedMath registry, supplied complete record with permanent ID \(30005694\), exact source problem code \texttt{catalogue entry}, title \emph{Projectivity and Cohomological Dimension without Boundedness or Unimodularity}. The record belongs to the Oberwolfach Reports Open Problems collection and names no individual proposer.
\item Supplied source citation: \emph{Mini-Workshop: Homological aspects of TDLC-groups} (2024), Oberwolfach Reports, pp.\ 3103--3138, DOI \texttt{10.4171/owr/2023/55}. \url{https://doi.org/10.4171/owr/2023/55}. The registry separately supplies proposed year 2023; the differing year designations are retained without reconciliation.
\item Supplied literature source, labelled ``Oberwolfach report'': \url{https://ems.press/content/serial-article-files/48166}.
\item Supplied literature source, attributed to Castellano--Marchionna--Weigel: \url{https://doi.org/10.1007/s00208-025-03116-7}. No article title or theorem statement is supplied for this reference.
\item Literature assessment dated exactly \texttt{2026-08-22}: the supplied curation labels the questions open, states that tensor-projectivity remains unanswered, and states that the 2025 Stallings--Swan result still assumes boundedness and unimodularity. This is supplied dated curation derived from the OWR record, not an independently verified literature assessment.
\item The registry labels the statement \texttt{corrected\_verified} and reports that it was checked against the official Oberwolfach report. These are supplied verification claims; the sources have not been independently consulted for this draft.
\end{itemize}
\end{document}
